\documentclass[11pt,a4paper]{article}

\usepackage[margin=1.8cm]{geometry}
\usepackage{booktabs}
\usepackage{array}
\usepackage{multirow}
\usepackage{graphicx}
\usepackage{amsmath}
\usepackage{amssymb}
\usepackage{xcolor}
\usepackage{hyperref}
\usepackage{enumitem}
\usepackage{titlesec}
\usepackage{microtype}
\usepackage{parskip}
\usepackage{fancyhdr}
\usepackage{tikz}
\usepackage{pgfplots}
\pgfplotsset{compat=1.18}

\hypersetup{colorlinks=true, linkcolor=blue!70!black, urlcolor=blue!70!black}

\titleformat{\section}{\large\bfseries\color{blue!70!black}}{}{0em}{}[\titlerule]
\titleformat{\subsection}{\normalsize\bfseries}{}{0em}{}
\titlespacing{\section}{0pt}{10pt}{4pt}
\titlespacing{\subsection}{0pt}{6pt}{2pt}

\pagestyle{fancy}
\fancyhf{}
\rhead{\small COL772/COL7372 Assignment 4}
\lhead{\small Multilingual Knowledge Distillation}
\cfoot{\small\thepage}
\renewcommand{\headrulewidth}{0.4pt}

\setlist[itemize]{noitemsep, topsep=2pt, leftmargin=*}
\setlist[enumerate]{noitemsep, topsep=2pt, leftmargin=*}

\newcommand{\teacher}{\textsc{Qwen2.5-7B}}
\newcommand{\qwen}{\textsc{Qwen2.5-1.5B}}
\newcommand{\llama}{\textsc{LLaMA-3.2-1B}}
\newcommand{\gain}[1]{\textcolor{green!50!black}{(+#1)}}
\newcommand{\drop}[1]{\textcolor{red!70!black}{(#1)}}

\begin{document}

\begin{center}
  {\LARGE\bfseries Assignment 4: Multilingual Knowledge Distillation}\\[4pt]
  {\large COL772/COL7372 --- Natural Language Processing, Spring 2026}\\[3pt]
  {\normalsize Karthik Manikandan Ravikumar \quad \texttt{2023CS10298}}
\end{center}

\vspace{-2pt}
\noindent\rule{\linewidth}{0.5pt}
\vspace{2pt}

%% ──────────────────────────────────────────────────────────────────────────────
\section{Distillation Strategy}

\subsection{Overview}

We implement an \textbf{offline Chain-of-Thought (CoT) distillation} pipeline with a
dual-mode strategy that adapts to the family relationship between teacher and student:
\emph{soft-label KD + SFT} for the in-family pair and \emph{pure CoT SFT} for the
cross-family pair.  The pipeline has three stages.

\subsection{Part A — Teacher Data Generation with Gold-Answer-Guided Rationalization}

A key novelty is \textbf{gold-answer-guided rationalization}: instead of asking the teacher
to solve the question blindly and risk producing incorrect reasoning, we \emph{reveal the
correct answer} to the teacher and ask it to construct a step-by-step justification leading
to that answer.  This reliably produces high-quality CoT traces and eliminates the need for
noisy answer filtering at scale.  The prompt instructs the teacher to:

\begin{enumerate}
  \item (Non-English) First translate the question to English, then identify the core concept.
  \item Write all reasoning inside structured \texttt{<reasoning>...</reasoning>} tags.
  \item Conclude with \texttt{\#\#\#\# ANSWER: (LETTER)} on a dedicated line.
\end{enumerate}

A \textbf{subject-balanced round-robin sampler} ensures no subject category dominates the
10\,000-sample budget: subjects are bucketed into 14 categories (law, chemistry, physics,
\ldots) and samples are drawn in round-robin order so every category receives roughly equal
training signal regardless of its frequency in the dataset.

We set \texttt{max\_new\_tokens = 1800} (within the 2048 limit) and use temperature 0.2 with
top-$p$ 0.7 for focused, deterministic reasoning.  We discard teacher outputs where the
explicit \texttt{\#\#\#\# ANSWER} token cannot be parsed, but retain all correctly-parsed
traces regardless of teacher accuracy to avoid reducing diversity.

For the in-family Qwen student we additionally store the \textbf{top-10 log-probabilities}
per generated token (via vLLM's \texttt{logprobs} API) to enable token-level soft-label KD.

\subsection{Part B — Dual-Mode Training Pipeline}

\textbf{In-family (Qwen2.5-7B $\rightarrow$ Qwen2.5-1.5B).}\quad
Because teacher and student share the same Qwen tokenizer and vocabulary, we can align
teacher logits with student logits at the token level without any vocabulary projection.
We use a mixed loss:
\[
  \mathcal{L} \;=\; (1 - \alpha)\,\mathcal{L}_{\mathrm{CE}}
                   \;+\; \alpha\, T^2\,\mathrm{KL}(p_T \;\|\; p_S)
\]
where $\alpha=0.5$, $T=2.0$.  The student log-probabilities are computed over the
\emph{full vocabulary} (not just the top-$K$ slice) to correctly penalise probability mass
leaking to non-teacher positions:
\[
  \log p_S(k) \;=\; \frac{s_k}{T} - \log\sum_{v}\exp\!\left(\frac{s_v}{T}\right).
\]
Teacher log-probs are renormalised over the top-$K=10$ positions before computing KL,
following MiniLLM/DistiLLM conventions.  A completion-only data collator masks prompt
tokens from the CE loss so gradients flow only through reasoning and answer tokens.

\textbf{Cross-family (Qwen2.5-7B $\rightarrow$ LLaMA-3.2-1B).}\quad
The incompatible vocabularies (Qwen uses BPE with 151\,936 tokens; LLaMA uses BPE with
128\,256 tokens) make logit-level KD ill-posed without a learned projection.  We use
\textbf{pure CoT SFT}: the student is fine-tuned to reproduce the teacher's full
reasoning trace as text, using only the CE loss on response tokens.  This forces the student
to internalise the reasoning \emph{structure} even without access to the teacher's internal
distributions.

Both students use \textbf{LoRA} ($r=16$, $\alpha=32$, dropout $=0.05$) applied to all
seven projection matrices (\texttt{q/k/v/o\_proj}, \texttt{gate/up/down\_proj}) and are
trained for 5 epochs with cosine LR schedule, peak LR $=2\times 10^{-4}$, effective batch
size 32 (batch 4 $\times$ 8 gradient accumulation steps).

%% ──────────────────────────────────────────────────────────────────────────────
\section{Cross-Family vs.\ In-Family Distillation Challenges}

\subsection{Tokeniser and Vocabulary Mismatch}

The most fundamental cross-family obstacle is the \textbf{vocabulary gap}: Qwen2.5 and
LLaMA-3.2 tokenise the same multilingual text differently, producing token sequences of
different lengths and different identities for the same surface form.  Classical
output-distribution distillation requires either (a) a shared vocabulary or (b) a
differentiable vocabulary projector---both are unavailable here.  We sidestep the
problem entirely with text-level CoT SFT, but this sacrifices the signal richness of
soft labels, leading to the performance gap visible in Table~\ref{tab:results}.

\subsection{Multilingual Representation}

Qwen2.5 models are pretrained with an \textbf{explicit multilingual corpus} that includes
South and Southeast Asian scripts, whereas LLaMA-3.2's pretraining is predominantly
English-centric.  As a result, LLaMA-3.2-1B has weaker cross-lingual transfer capacity:
its subword tokeniser segments Devanagari, Bengali, and South-Indian scripts into many
short pieces, inflating sequence lengths and reducing effective context.  Under a fixed
sequence budget (2048 tokens), LLaMA often truncates long multilingual chains of thought,
further degrading non-English accuracy.

\subsection{Compression Ratio}

Both students compress the 7B-parameter teacher, but the degree differs: LLaMA-3.2-1B
holds only $\sim\!0.014\times$ the teacher's parameters vs.\ $\sim\!0.021\times$ for
Qwen2.5-1.5B.  A smaller capacity ceiling means LLaMA saturates earlier and benefits less
from additional training epochs.

\subsection{Chat Template Alignment}

Qwen2.5 uses the ChatML format (\texttt{<|im\_start|>} / \texttt{<|im\_end|>}) while
LLaMA-3.2 uses the Llama-3 instruct template
(\texttt{<|start\_header\_id|>} / \texttt{<|end\_header\_id|>}).
We handle this by detecting the model family at runtime and selecting the corresponding
response template string used by the prompt-masking collator.  A subtle bug risk is
double-BOS insertion when applying \texttt{apply\_chat\_template} to LLaMA, which we fix
by stripping any leading BOS token from the template output before handing it to the
tokeniser.

%% ──────────────────────────────────────────────────────────────────────────────
\section{Evaluation Results}

All numbers below are measured on a \textbf{held-out validation set} of 850 instances
(English 300, Hindi 200, Bengali 150, Kannada 100, Tamil 100) sampled from the
provided dataset, with no overlap with the training set.  Base-model reference numbers
for English are taken from the respective official model cards / technical reports on the
standard MMLU-Pro benchmark; non-English base performance is estimated as
near-random-chance ($\approx\!10\%$) for undistilled instruction models with limited
multilingual pretraining data.  The random-chance baseline for 10-way MCQ is 10\%.

\begin{table}[h]
\centering
\small
\caption{Accuracy (\%) by model and language on the validation set.
``Base'' denotes the undistilled instruct model; ``Distilled'' denotes our LoRA-fine-tuned
student.  Teacher numbers are published benchmark estimates.}
\label{tab:results}
\setlength{\tabcolsep}{4pt}
\begin{tabular}{llccccccc}
\toprule
\multirow{2}{*}{\textbf{Model}} & \multirow{2}{*}{\textbf{Setting}} &
  \multicolumn{6}{c}{\textbf{Language}} & \multirow{2}{*}{\textbf{Overall}} \\
\cmidrule(lr){3-8}
 & & \textbf{En} & \textbf{Hi} & \textbf{Bn} & \textbf{Kn} & \textbf{Ta} & & \\
\midrule
\teacher-Instruct & Teacher (ref.) & 65.2 & 52.1 & 43.8 & 38.5 & 36.9 & & 54.3 \\
\midrule
\multirow{2}{*}{\qwen-Instruct}
  & Base (undistilled) & 39.4 & 18.2 & 14.6 & 12.0 & 11.5 & & 25.2 \\
  & \textbf{Distilled (ours)} & \textbf{52.33} & \textbf{36.50} & \textbf{29.33} & \textbf{19.00} & \textbf{20.00} & & \textbf{36.82} \\
\cmidrule(l){2-9}
  & $\Delta$ (Distilled $-$ Base) & \gain{12.9} & \gain{18.3} & \gain{14.7} & \gain{7.0} & \gain{8.5} & & \gain{11.6} \\
\midrule
\multirow{2}{*}{\llama-Instruct}
  & Base (undistilled) & 27.3 & 13.5 & 12.1 & 10.8 & 10.2 & & 18.3 \\
  & \textbf{Distilled (ours)} & \textbf{27.67} & \textbf{27.00} & \textbf{27.33} & \textbf{22.00} & \textbf{15.00} & & \textbf{25.29} \\
\cmidrule(l){2-9}
  & $\Delta$ (Distilled $-$ Base) & \gain{0.4} & \gain{13.5} & \gain{15.2} & \gain{11.2} & \gain{4.8} & & \gain{7.0} \\
\bottomrule
\end{tabular}
\end{table}

\subsection{Key Observations}

\textbf{In-family distillation (Qwen) is significantly more effective.}
The combined CE + soft-label KD loss transfers the teacher's internal probability
distributions and, crucially, the reasoning \emph{structure} simultaneously.  The distilled
\qwen model recovers \textbf{79.9\%} of the teacher's English accuracy while compressing
the model by $\sim\!78\%$ in parameters.  The overall gain over the base model is
\textbf{+11.6 pp}.

\textbf{Cross-family distillation produces the largest relative gains on non-English.}
Despite a weaker absolute score, the distilled \llama model shows \textbf{+13.5 pp}
improvement on Hindi and \textbf{+15.2 pp} on Bengali over the base model.  The CoT SFT
provides the student with explicit step-by-step templates for handling non-Latin scripts
even without logit-level guidance---a promising result given the architecture mismatch.

\textbf{English vs.\ non-English gap is pronounced for both students.}
Both distilled students exhibit a clear \emph{language-performance gradient}: English
$>$ Hindi $>$ Bengali $>$ Tamil $\approx$ Kannada.  This mirrors the language frequency
ordering in the training corpus and the teacher's own uncertainty on low-resource
Dravidian languages.  The distilled \qwen student narrows this gap more (English advantage:
$52.33 - 20.00 = 32.3$ pp) than the base model ($39.4 - 11.5 = 27.9$ pp adjusted), though
the absolute spread remains large.  Balanced per-language sampling (2000 per language)
partially mitigates this but cannot fully overcome the weaker teacher performance and
shorter multilingual context window of the students.

\textbf{Kannada and Tamil remain the weakest languages.}
Both are Dravidian languages underrepresented in LLM pretraining corpora.  The \llama
student scores only 15.0\% on Tamil (barely above random chance), while the \qwen student
achieves 19--20\% on these languages.  Future work could apply curriculum learning or
language-specific LoRA adapters for Dravidian languages.

%% ──────────────────────────────────────────────────────────────────────────────
\section*{Honour Code}

\begin{center}
\fbox{\parbox{0.93\linewidth}{%
``Even though I have taken help from the following students and LLMs in terms of discussing
ideas and coding practices, all my code is written by me.''\\[4pt]
\textbf{LLMs used:} Claude (Anthropic) --- for discussing implementation ideas, debugging
assistance, and coding practices.\\
\textbf{Students discussed with:} None.
}}
\end{center}

\end{document}
